\section{Conclusion}\label{sec:conclusion}
This paper studied when a multi-UAV replanner should stop computing and publish its plan. AnytimeSwarm stops an anytime
ALNS when a learned prediction of the next improvement no longer exceeds a delay cost derived from the slack of the
plan, and AnytimeSwarm-R extends this cost with a calibrated model of execution errors and latency. Both rules were
evaluated in a pre-registered study with a kinematic tier and a PX4 software-in-the-loop tier on held-out missions.

With exact execution, AnytimeSwarm reached a lower regret than fixed budgets and stagnation stopping at equal
computation on every fleet scale. With uncertain execution, the adaptive rules served as much reward as the best fixed
budget, which had been matched to them on development data, and avoided the losses of long budgets on slow planners
and of short budgets on large fleets without tuning per scale. AnytimeSwarm-R needed \val{JMsRedAS}\% less computation
than AnytimeSwarm in the kinematic tier and \val{JCMsRedAS}\% less in the PX4 tier. Its gain in utility was small and not significant, so the execution-aware delay cost is
best seen as a way to stop earlier at no loss rather than as a way to serve more reward. Periodic or deviation-triggered
replanning raised the utility of every method more than any choice of stopping rule. The kinematic tier predicted the
utility of PX4 flights closely, which supports its use for large comparisons, but differences of less than one
percentage point between methods remained within the noise of both tiers.

Future work will let the error model shape the routes themselves through chance-constrained planning, choose the
replanning trigger and the stopping rule jointly, and run the planner on embedded hardware and in field tests, where the
reduction of computation is expected to matter most.
